\section{Design}
\label{sec:design}

\subsection{The programming model}

\modelly{} is an independent implementation: it contains no Mesa code and does not depend on
Mesa. Its compatibility with Mesa is at the level of the programming model and the API's
names, which is what lets a codemod migrate a Mesa model (\cref{sec:samecode}).
A \modelly{} model is a \code{Model} subclass that owns its agents, its spaces, its clock and
its random generators; an agent is an \code{Agent} subclass with a \code{step} method or
whatever behaviour the model calls. Parameters live in a typed \code{Scenario} that also
carries the seed, so that a model's whole input is one object that can be stored, swept and
hashed. Agents are kept in ordered \code{AgentSet}s with \code{do}, \code{shuffle\_do},
\code{select}, \code{agg} and \code{groupby}; the model's registries are backed by
dictionaries, so membership tests run in C and iteration order is deterministic. Time advances
in ticks, through scheduled events, or both on one heap with lazy cancellation.

\subsection{Three ways to store a population}
\label{sec:storage}

The same model can be written at three levels of the speed/expressiveness trade-off, and
\modelly{} supports all three in one framework:

\begin{description}[style=nextline,leftmargin=1.5em,itemsep=2pt]
  \item[Objects (the per-agent style of Mesa).] Each agent is a Python object with its own attributes, and each
  cell keeps the list of its occupants. This is the most direct translation of a model's
  description and the easiest to read.
  \item[Columns.] An attribute declared as a \code{Column} is stored in a model-owned
  \code{numpy} array, one table per agent class. Each agent still reads and writes it as an
  attribute (a property reading its row), while the whole population is available as one array
  for vectorised phases (\cref{lst:columns}). Phases that act on all agents at once (movement,
  infection, ageing, mortality) become array operations.
  \item[Cohorts.] When many agents are identical in every state variable that matters (stage,
  age, sex), a \code{CohortTable} stores one row per class with a count, and stochastic events
  become binomial draws on the counts (below), the super-individual approach of
  individual-based ecology \citep{scheffer1995super}. A population of tens of thousands then
  costs at most about 190 rows in our runs.
\end{description}

\paragraph{When cohorts keep the individual model's law.} The claim is conditional, and
we state its conditions. Let a row hold $n$ individuals in state $x$, and let each face,
independently, the same probability $s(x)$ of surviving a phase. In the individual model
the survivors are $\sum_{k=1}^{n} B_k$ with $B_k \sim \mathrm{Bernoulli}(s(x))$
independent, which is $\mathrm{Binomial}(n, s(x))$; the cohort draws that count directly:
\begin{equation}
  \begin{aligned}
    &\text{survival (\code{thin}):} && n' \sim \mathrm{Binomial}\bigl(n, s(x)\bigr),\\
    &\text{transition (\code{split}):} && K \sim \mathrm{Binomial}\bigl(n, p(x)\bigr)
      \text{ move to state } x', \text{ and } n - K \text{ stay},\\
    &\text{events (\code{draw}):} && E \sim \mathrm{Binomial}\bigl(n, q(x)\bigr).
  \end{aligned}
  \label{eq:cohorts}
\end{equation}
The three operations cover survival, two-way transitions and counts of events that each
individual triggers at most once per phase (an egg batch, an emergence).

\begin{description}[style=nextline,leftmargin=1.5em,itemsep=2pt]
  \item[Proposition.] Suppose that, in every phase, (A1) all the individuals of the
  population, within a row and across rows, are conditionally independent given the state
  vector at the start of the phase; (A2) every individual of a row faces the same
  probabilities, which are functions of that state vector only (homogeneity within the row
  and a synchronous update); and (A3) each individual's fate in the phase is one of the
  three operations above (one draw per individual, into at most two outcomes). Then, by
  induction over the phases, the vector of counts per state has the same distribution in
  the cohort model as in the individual model.
\end{description}

\noindent The proof is the binomial identity above applied row by row. Conditional on the
state vector at the start of the phase, (A1) makes the fates of all individuals
independent, so the fates of distinct rows are independent of one another and, within a
row, (A2) and (A3) make them exchangeable Bernoulli variables with the row's probability:
each row's count is therefore binomial, the rows' counts are conditionally independent,
and their joint law given the state is the one the cohort model draws from (one
independent binomial per row, with the same probabilities, since the cohort evaluates them
from the same state vector). The induction over phases carries the equality of the joint
law of the state vector from one phase to the next. Independence across rows is part of
(A1), not a consequence of independence within rows: a phase in which two rows' fates are
coupled (a shared resource drawn once and allocated to both, say) breaks (A1) even when
each row's members are independent among themselves. The proposition says nothing when an
assumption fails, and two cases show that the conclusion then fails too
(20,000 replicates of one phase on $n = 200$ individuals with $s = 0.8$, by
\code{paper/scripts/cohort\_counterexample.py}, whose data file is next to it). First, an
asynchronous update breaks (A2): when each individual's survival probability depends on the
count alive at its own turn (rising by 0.002 per earlier death), the individual model keeps
167.0 survivors on average and the cohort, which evaluates the rate once at the start of the
phase, 160.0, a difference of 7.1 (standard error 0.05). Second, hidden heterogeneity breaks
(A2) in another way: when the individuals of a row carry a variable the row does not (each
its own probability, uniform on $0.8 \pm 0.2$), the survivors are Poisson-binomial, with the
same mean as the cohort's binomial (161.0 against 160.9) but a smaller variance (28.8 against
31.3; exact values 28.8 and 31.4; median-centred Levene $p = 10^{-6}$). Interactions within
a phase between members of the same row (a pairwise contest, say) break (A1) in the same
way. In the positive case, where the assumptions hold, the two distributions coincide (means
160.0 and 160.0, variances 31.9 and 32.2 for an exact 32.0, Kolmogorov--Smirnov $p = 0.997$).
The practical rule is therefore: a cohort model reproduces the individual model that updates
synchronously with rates evaluated once per phase, as the Mesa implementations of the insect
playbook do (\cref{sec:agreement} checks the insect model's statistics), and any state
variable that changes an individual's probabilities must be part of the row's key.
Individual identities, and histories that differ within a state, are not represented; a
model that needs them stays at the object or column level.

\begin{lstlisting}[float=tp,caption={A wealth-exchange model whose agents keep their wealth in a column;
the step acts on the whole population at once. Run as shown with \modelly{} 1.4.2.},
label={lst:columns}]
import numpy as np
from modelly import Agent, Column, Model
from modelly.scenarios import Scenario

class Economy(Scenario):
    n: int = 1000

class Trader(Agent):
    wealth = Column(int, default=1)

class Exchange(Model):
    def __init__(self, scenario: Economy | None = None):
        scenario = scenario or Economy()
        super().__init__(scenario=scenario)
        Trader.create_agents(self, scenario.n)

    def step(self):
        wealth = self.columns[Trader].wealth   # the whole population, one array
        givers = np.flatnonzero(wealth > 0)
        takers = self.rng.integers(0, wealth.size, size=givers.size)
        np.subtract.at(wealth, givers, 1)
        np.add.at(wealth, takers, 1)

model = Exchange(Economy(n=1000, rng=42))
model.run_for(100)       # same seed, same run
\end{lstlisting}

The columns view and the per-agent attribute are the same memory, so a model can mix
per-agent code and vectorised phases freely: the benchmarks in \cref{sec:performance} include
variants of the same model at each level.

\subsection{Seeded randomness}

Every model owns two generators: \code{model.random}, a \code{FastRandom} that produces exactly
the numbers of Python's \code{random.Random} for the same seed (the Mersenne Twister
\citep{matsumoto1998mersenne}), and \code{model.rng}, a \code{numpy} \code{Generator}.
\code{FastRandom}'s hot methods (\code{randint}, \code{choice}, \code{shuffle},
\code{sample}) run in C on CPython's own generator state; at import it is checked against
\code{random.Random}, and a pure-Python implementation with the same stream is used when the C
module is absent. A model migrated from code that used \code{random.Random(seed)} therefore keeps
its exact stream, which is what makes the bit-identical results of \cref{sec:samecode}
possible. A model can instead give every agent generators of its own
(\code{agent\_streams = True}): each is seeded from the scenario's seed sequence and the
agent's stream key, so what an agent draws no longer depends on the activation order or on
the population, which is what lets a run split across workers reproduce the single-process
run (\cref{sec:split}). The flag is off by default; turning it on changes a model's numbers.

\subsection{Compiled loops with pure-Python twins}

Five small C modules implement the framework's inner loops: the random core; per-agent method
calls (\code{do}, \code{shuffle\_do}) and the breadth-first neighbourhood walk; the cell move;
set filtering and attribute gathering (\code{select}, \code{agg}); and the event list's heap (\cref{app:perf2}). Each C function is a
line-for-line translation of a \modelly{} Python function that stays in the source as its
fallback; tests run the
two side by side, and \code{MODELLY\_PURE\_PYTHON=1} forces the Python path. Results are
identical either way. The modules build with GCC, Clang and MSVC; the release workflow builds
them on the tag for Linux (glibc and musl, x86-64 and ARM64), macOS and Windows (x86-64 and
ARM64), where a module that fails to compile fails the
build and each wheel's test imports four of the five compiled (all but the event heap), checked
as \cref{tab:determinism} states. They also build for free-threaded CPython,
where critical sections and cache generation counters replace the global interpreter lock,
and the compiled modules' tests, except the event heap's, run on that build on the development
machine.

\subsection{Spaces}

Discrete spaces (orthogonal Moore and von Neumann grids, hexagonal grids, networks, Voronoi
meshes) share one cell abstraction with cached neighbourhoods, property layers backed by arrays
and an empty-cell index; a \code{CellIndex} exposes any discrete space as integer rows and a
sparse adjacency for vectorised counts and random-neighbour draws. Continuous space keeps
positions in one array and answers neighbour queries with one KD-tree per step. A geographic
space, comparable in scope to Mesa-Geo \citep{wang2022mesageo}, projects coordinates once and supports regions, contiguity and road networks, and
temporal and multilayer networks represent contact structures that change over time.

\subsection{State that survives a pause}

\code{model.save(path)} writes a checkpoint holding the whole model, both generators and the
event scheduler; a run resumed from it, in a new process, is the run it would have been.
Agents are written flat, so the cost of a checkpoint does not grow with the depth of the
agent--cell graph. Checkpoints are the basis of time travel, resumable experiments and
calibrations, and Studio's timeline (\cref{sec:workflow}).
